\subsubsubsection{So sánh các kiến trúc hiện có}

\clearpage
\begin{landscape}
    \begin{table}[h!]
        \centering
        \setlength{\tabcolsep}{4pt}
        \footnotesize
        \begin{tabularx}{\linewidth}{|l|X|X|X|X|}
            \hline
            \textbf{Tiêu chí}                       & \textbf{Layered Architecture}                                                                                                                                     & \textbf{Modular Monolith}                                                                                                                                                                                                                   & \textbf{Microservices Architecture}                                                                                                                             & \textbf{Event-Driven Architecture}                                                                                                                                          \\
            \hline
            \textbf{Độ phức tạp phát triển}         & \textbf{Thấp.} Cấu trúc đơn giản, tuyến tính và quen thuộc. Toàn bộ mã nguồn nằm trong một dự án, dễ dàng điều hướng và phát triển ban đầu.                       & \textbf{Thấp.} Tổ chức theo domain giúp code có cấu trúc rõ ràng hơn, nhưng cần thiết kế ranh giới module cẩn thận ngay từ đầu. Vẫn chạy trong một tiến trình duy nhất nên không phát sinh độ phức tạp của hệ thống phân tán.               & \textbf{Cao.} Là một hệ thống phân tán, đòi hỏi xử lý các vấn đề phức tạp như giao tiếp mạng giữa các dịch vụ, service discovery, và quản lý cấu hình phân tán. & \textbf{Rất cao.} Vừa là hệ thống phân tán, vừa hoạt động bất đồng bộ. Lập trình viên phải tư duy theo luồng sự kiện, xử lý các vấn đề như thứ tự sự kiện, xử lý trùng lặp. \\
            \hline
            \textbf{Khả năng mở rộng}               & \textbf{Thấp.} Phải mở rộng toàn bộ ứng dụng theo chiều ngang. Kém hiệu quả khi chỉ một phần của hệ thống bị quá tải.                                             & \textbf{Trung bình.} Tương tự Layered, toàn bộ ứng dụng vẫn được triển khai và mở rộng như một đơn vị duy nhất. Tuy nhiên, ranh giới module rõ ràng tạo nền tảng thuận lợi để tách thành Microservices trong tương lai nếu cần.             & \textbf{Rất cao.} Cho phép mở rộng từng dịch vụ một cách độc lập và linh hoạt. Tối ưu hóa việc sử dụng tài nguyên phần cứng.                                    & \textbf{Rất cao.} Các thành phần sản xuất và tiêu thụ sự kiện có thể được mở rộng độc lập.                                                                                  \\
            \hline
            \textbf{Khả năng bảo trì \& sửa đổi}    & \textbf{Cao khi dự án nhỏ, thấp khi lớn dần.} Khi mã nguồn phình to, sự liên kết chặt chẽ giữa các module làm cho việc sửa đổi một nơi có thể gây lỗi ở nơi khác. & \textbf{Cao.} Mỗi module có ranh giới trách nhiệm rõ ràng, giúp dễ dàng tìm kiếm và sửa đổi code. Sự thay đổi trong một module ít có khả năng gây ra hiệu ứng lan truyền sang các module khác.                                              & \textbf{Cao.} Các dịch vụ nhỏ, độc lập và có phạm vi rõ ràng nên rất dễ hiểu, dễ sửa đổi và nâng cấp mà không ảnh hưởng đến toàn bộ hệ thống.                   & \textbf{Trung bình đến cao.} Các khớp nối lỏng giúp việc sửa đổi các dịch vụ dễ dàng. Tuy nhiên, việc hiểu và theo dõi toàn bộ luồng nghiệp vụ qua các sự kiện sẽ phức tạp. \\
            \hline
            \textbf{Triển khai}                     & \textbf{Đơn giản.} Chỉ cần triển khai một đơn vị duy nhất. Quy trình CI/CD đơn giản.                                                                              & \textbf{Đơn giản.} Giống Layered Architecture, toàn bộ ứng dụng được đóng gói và triển khai như một artifact duy nhất. Quy trình CI/CD đơn giản, phù hợp với nhóm nhỏ.                                                                      & \textbf{Phức tạp.} Phải triển khai và quản lý vòng đời của nhiều dịch vụ. Yêu cầu các công cụ điều phối như Kubernetes và một quy trình DevOps hoàn chỉnh.      & \textbf{Phức tạp.} Tương tự Microservices, nhưng cần thêm việc triển khai và quản lý hệ thống message broker.                                                               \\
            \hline
            \textbf{Khả năng phục hồi \& chịu lỗi}  & \textbf{Thấp.} Lỗi nghiêm trọng ở một module có thể làm sập toàn bộ ứng dụng.                                                                                     & \textbf{Thấp.} Lỗi nghiêm trọng vẫn có thể ảnh hưởng toàn bộ hệ thống. Tuy nhiên, sự phân tách module rõ ràng giúp việc cô lập và xử lý lỗi dễ dàng hơn.                                                 & \textbf{Cao.} Lỗi ở một dịch vụ (không phải dịch vụ cốt lõi) sẽ không làm ảnh hưởng đến các dịch vụ khác.                                                       & \textbf{Cao.} Nếu một dịch vụ bị lỗi, sự kiện có thể được lưu lại và xử lý sau.            \\
            \hline
        \end{tabularx}
        \caption{So sánh các kiến trúc hiện có}
        \label{tab:architecture_comparison}
    \end{table}
\end{landscape}

Dựa trên bảng phân tích chi tiết theo từng tiêu chí và đối chiếu với bối cảnh cụ thể của dự án ``Gợi ý hành trình du lịch dựa trên AI và đánh giá cộng đồng'', nhóm quyết định lựa chọn \textbf{Modular Monolith} vì những lý do sau:
\begin{itemize}
    \item \textbf{Độ phức tạp phù hợp:} Với quy mô của một dự án luận văn tốt nghiệp, được phát triển bởi một đội ngũ có giới hạn, Modular Monolith cung cấp sự cân bằng lý tưởng. Kiến trúc này kế thừa sự đơn giản trong triển khai của Layered Architecture (một artifact duy nhất, quy trình CI/CD đơn giản), trong khi tổ chức mã nguồn theo domain giúp tránh tình trạng ``Big Ball of Mud'' khi hệ thống mở rộng với nhiều tính năng.
    \item \textbf{Phù hợp với độ phức tạp của nghiệp vụ:} Hệ thống TravelEZ có nhiều domain rõ ràng và tách biệt: quản lý điểm tham quan (POI), lộ trình, cộng đồng (bài đăng, bình luận), nhắn tin, thông báo, và tích hợp AI. Mỗi domain được đóng gói thành một module độc lập với ranh giới trách nhiệm rõ ràng, giúp việc phát triển song song giữa các thành viên trong nhóm trở nên thuận lợi hơn.
    \item \textbf{Nền tảng cho việc phát triển trong tương lai:} Ranh giới module rõ ràng tạo tiền đề để tách từng module thành Microservices riêng biệt khi hệ thống có nhu cầu mở rộng quy mô trong tương lai, mà không cần phải tái cấu trúc toàn bộ mã nguồn.
    \item \textbf{Tránh sự phức tạp không cần thiết:} Việc áp dụng Microservices hay Event-Driven Architecture cho dự án này sẽ là một trường hợp ``over-engineering''. Lợi ích mà chúng mang lại không đủ để bù đắp cho chi phí về thời gian, công sức và độ phức tạp về hạ tầng mà đội ngũ phải bỏ ra để quản lý chúng.
\end{itemize}
